\subsection{Fréchet 空间对偶上的紧收敛拓扑}

定理 1（Banach-Dieudonné）. --- 设 E 是一个局部凸可度量化空间。下列拓扑在 E 的对偶 E' 上一致：

\begin{enumerate}
\def\labelenumi{\alph{enumi})}
\item
  拓扑 \(\mathcal{T}_{\mathfrak{N}}\)（\(\mathfrak{N}\)-收敛的），其中 \(\mathfrak{N}\) 是 \(E\) 的这样的子集组成的族：每个子集由一个收敛于 0 的序列的点组成；
\item
  在 E 的紧子集上一致收敛的拓扑 \(\mathcal{T}_{\mathrm{c}}\)；
\item
  在 E 的预紧子集上一致收敛的拓扑 \(\mathcal{T}_{pc}\)；
\item
  拓扑 \(\mathcal{T}_f\)，它是这样的最细拓扑：在每个等度连续子集上诱导出与 \(\sigma(\mathrm{E}',\mathrm{E})\) 相同的拓扑（这些子集属于 \(\mathrm{E}'\)）。
\end{enumerate}

首先注意到，\(E'\) 的子集 A 对于 \(T_{f}\) 闭，当且仅当 \(A \cap H\) 对于 \(\sigma(E', E)\) 闭（对每个子集 H，它属于 \(E'\)、等度连续且对于 \(\sigma(E', E)\) 闭）。弱拓扑 \(\sigma(E', E)\) 与 \(T_{pc}\) 在 \(E'\) 的每个等度连续子集上诱导出相同的拓扑（III，第 17 页，命题 5）。因此，拓扑 \(T_{R}\)、\(T_{c}\)、\(T_{pc}\)、\(T_{f}\) 中的每一个都比它后面一个粗。于是只需证明 \(T_{R}\) 比 \(T_{f}\) 细。此外，\(E'\) 中的每个平移对于 \(T_{f}\) 是一个同胚。因此只需证明：若 F 是 \(E'\) 的一个对于 \(T_{f}\) 闭且不包含 0 的子集，则存在集 \(S \in R\) 使得 \(S^{\circ} \cap F = \varnothing\)。

设 \((\mathbf{U}_{n})_{n\geqslant0}\) 是 E 中零邻域的一个递减序列，构成零邻域的一个基本系。我们将对 \(n\geqslant0\) 用归纳法构造有限集 \(X_{n}\)，使得我们有

\[
\mathrm{X} _ {n} \subset \mathrm{U} _ {n}\tag{4}
\]

\[
\left(\bigcup_ {0 \leqslant p \leqslant n} \mathrm{X} _ {p}\right) ^ {\circ} \cap \mathrm{U} _ {n + 1} ^ {\circ} \cap \mathrm{F} = \varnothing\tag{5}
\]

对每个整数 \(n \geqslant 0\)。设 \(m \geqslant 0\) 是这样的整数：\(X_{n}\) 已对 \(0 \leqslant n < m\) 构造出来，且对 \(0 \leqslant n < m\) 满足 (4) 与 (5)。对每个 \(x \in U_{m}\)，令

\[
\mathrm{F} _ {x} = \big (\bigcup_ {0 \leqslant p <   m} \mathrm{X} _ {p} \big) ^ {\circ} \cap \{x \} ^ {\circ} \cap \mathrm{U} _ {m + 1} ^ {\circ} \cap \mathrm{F}.
\]

取 \(n = m - 1\) 的公式 (5) 蕴含 \(\bigcap_{x\in \mathrm{U}_m}\mathrm{F}_x = \emptyset\)。此外，集 \(\mathrm{U}_{m + 1}^{\circ}\) 是等度连续的，且对于 \(\sigma (\mathrm{E}',\mathrm{E})\) 紧。鉴于 \(\mathcal{T}_f\) 的定义，每个集 \(\mathrm{F}_x\) 对于 \(\sigma (\mathrm{E}',\mathrm{E})\) 紧；因此存在有限子集 \(\mathrm{X}_m\)（\(\mathrm{U}_m\) 中的）使得 \(\bigcap_{x\in \mathrm{U}_m}\mathrm{F}_x = \emptyset\)，即关系 (5) 对 \(n = m\) 满足。

令 \(S = \bigcup_{n\geqslant 0}X_n\)。我们有 \(X_{n}\subset U_{p}\)（对 \(n\geqslant p\)），因此 \(S\) 是 \(E\) 中一个收敛于 0 的序列的点组成的集。由 (5) 我们推出 \(S^{\circ}\cap U_{n + 1}^{\circ}\cap F = \emptyset\)，且由于 \(E'\) 是诸集 \(U_{n + 1}^{\circ}\) 组成的序列之并，我们得到 \(S^{\circ}\cap F = \emptyset\)。

推论 1. --- 设 E 是一个局部凸可度量化空间。E 的每个预紧子集都包含在一个收敛于 0 的序列的点组成的集的闭凸均衡包中。

这由拓扑 \(\mathcal{T}_{pc}\) 与 \(\mathcal{T}_{\mathfrak{N}}\) 相同这一事实得出，依据是 III 第 15 页的命题 2。

推论 2. --- 设 E 是一个 Fréchet 空间。为使 E 的对偶 E' 的凸子集 A 对于 \(\sigma(E', E)\) 闭，必要且充分的是 \(A \cap U^{\circ}\) 对于 \(\sigma(E', E)\) 闭（对 E 中 0 的每个邻域 U）。

由于 E 完备，拓扑 \(T_{c}\) 在 \(E'\) 上与 \(E'\) 和 E 之间的对偶相容（IV，第 3 页，例）；因此 \(E'\) 中的闭凸子集对于 \(T_{c}\) 与 \(\sigma(E', E)\) 是相同的（IV，第 1 页，命题 1）。于是推论由拓扑 \(T_{c}\) 与 \(T_{f}\) 的相同性得出。

回忆（I，第 13 页）：E' 中对于 \(\sigma(\mathrm{E}', \mathrm{E})\) 闭的超平面是 E' 上由 E 的元素所对应的线性型的核。因此推论 2 给出了 III 第 21 页的推论 1（对 Fréchet 空间而言）的另一个证明。

推论 3. --- 设 E 是一个 Banach 空间，M 是 E 的对偶 E' 的一个向量子空间。为使 M 对于弱拓扑 \(\sigma(E', E)\) 闭，必要且充分的是它与 E' 中（闭）单位球的交对于 \(\sigma(E', E)\) 闭。

例. --- * 设 H 是一个满足第一可数公理的 hilbertian 空间；设 \(H_{\sigma}\) 表示赋以弱化拓扑的空间 H。设 \(\mathcal{L}^{1}(H)\) 是 H 的核型自同态组成的 Banach 空间（V，第 51 页，及 TS，V）；\(\mathcal{L}^{1}(H)\) 中的范数由 \(\|u\|_{1} = \operatorname{Tr}((u^{*}u)^{1/2})\) 定义。我们可以把 \(\mathcal{L}(H)\) 等同于 Banach 空间 \(\mathcal{L}^{1}(H)\) 的对偶，方法是把线性型 \(\phi_{u}: v \mapsto \operatorname{Tr}(uv)\)（\(\mathcal{L}^{1}(H)\) 上的）与每个 \(u \in \mathcal{L}(H)\) 相对应。设 A 是 \(\mathcal{L}(H)\) 的一个子代数，包含 1 且在 \(u \mapsto u^{*}\) 下稳定；它是 von Neumann 代数当且仅当它在 \(\mathcal{L}(H)\) 中对于弱拓扑 \(\sigma(\mathcal{L}(H), \mathcal{L}^{1}(H))\) 闭。由推论 3，我们推出下列判别法：为使 A 是 von Neumann 代数，必要且充分的是：若 \((u_{n})\) 是 A 中范数 \(\leqslant 1\) 的元素的任一序列，在空间 \(\mathcal{L}_{s}(H; H_{\sigma})\) 中有极限 u，则 u 属于 A。

